
\textbf{7.16.2.}
Nous ne passerons pas en revue ici, dans le cas des fibres géométriques, l'extension de toutes les propriétés des homomorphismes de spécialisation prouvées dans les sections précédentes pour les fibres «~arithmétiques~»~; elles se déduisent toutes trivialement de ces dernières par passage à la limite. Signalons seulement la conséquence intéressante suivante de 7.9.4, lorsque @@M0@@ est lisse et propre : @@M1@@ est isomorphe à un quotient d'un sous-groupe de @@M2@@, en particulier son rang est majoré par celui de ce dernier groupe. C'est donc un résultat de \emph{semi-continuité} pour le rang du groupe de cycles à équivalence numérique près.

\subsection*{7.17. Cas d'une base quelconque.}
Soit maintenant
@@M3@@
un morphisme plat et quasi-compact de préschémas noethériens, soient @@M4@@ deux points de @@M5@@ tels que @@M6@@ soit spécialisation de @@M7@@, on veut mettre en relation les groupes @@M8@@ etc pour les deux fibres @@M9@@ et @@M10@@, et pour des fibres géométriques correspondantes. Pour ceci, on note (EGA II 7.1.9) qu'il existe un morphisme
@@M11@@
où @@M12@@ est un trait, et où @@M13@@ transforme le point fermé @@M14@@ en @@M15@@, le point générique @@M16@@ en @@M17@@~; on peut même supposer que @@M18@@. Supposons que @@M19@@ soit noethérien, ce qui est le cas si @@M20@@ est de type fini ou si @@M21@@ est essentiellement de type fini sur @@M22@@ (EGA IV 1.3.8) (ce qu'on peut supposer si @@M23@@ lui-même est universellement japonais (EGA @@M24@@ 23.1.1)). Alors l'homomorphisme de spécialisation pour @@M25@@ définit un homomorphisme
@@M26@@
